\subsection{Short even lists: the fourth moment and all outer lengths}\label{proof:R2-C041}

\paragraph{R41.1 Exact theorem and actual integer direction table}

All expectations and operators are on the unchanged original finite probability space. Let \(K\) be real skew, pointwise bounded by one, and put

\fitdisplay{
B=-K,\quad u=B\one,\quad T=J+B,\quad A=J-B=T^*,\quad
 p=1-u,\quad q=1+u,
}{R41.L60}

\fitdisplay{
\begin{gathered}
f_n=A^n\one,\quad \bar f_n=T^n\one,\quad h_n=Tf_{n-1},\quad W_n=A\bar f_{n-1},\\
\qquad \delta_i=\|B^i\one\|^2,\quad d=\delta_1,\quad\theta=13/32.
\end{gathered}
}{R41.L64}

The terminal has an actual strict even residual

\fitdisplay{
\begin{gathered}
a_1<b_1<\cdots<a_j<b_j,\qquad j\ge2,\quad a_1\in\{2,3\},\\
\qquad F=\prod_{i=1}^j h_{a_i}f_{b_i}.
\end{gathered}
}{R41.L69}

There is an existing original anchor \(m\ge3\) at the receiving vertex. Optional original negative posts and genuine equal pairs are

\fitdisplay{
\begin{gathered}
P=\prod_{i=1}^k W_{u_i}\bar f_{v_i},\quad\\
4\le u_1<v_1<\cdots<u_k<v_k,\quad k\ge0;\qquad E_r,E_t.
\end{gathered}
}{R41.L75}

The blocks may overlap and be dependent.

Choose the actual anchor \(X\) and actual connector \(S_L\) as follows.

\begin{description}[style=nextline,leftmargin=1.6em,before=\raggedright]
\item[{$a_1=2$, $L=1$}] Use \(X S_L=\bar f_m A\).
\item[{$a_1=2$, $L\ge2$}] Use \(X S_L=W_mT^L\).
\item[{$a_1=3$, $m\le L$}] Use \(X S_L=W_mT^L\).
\item[{$a_1=3$, $m>L$}] Use \(X S_L=\bar f_mAT^{L-1}\).
\end{description}

Set

\fitdisplay{
H_D=E_r P X S_L(E_t F).
}{R41.1.1}

Every direction is fixed solely from original integers. In an \(A T^{L-1}\) row only the receiver-adjacent connector edge is A; all remaining connector edges are \(T\). The \(b,\allowbreak \allowbreak  W,\allowbreak \allowbreak  h\) and \(f\) factors are the actual paths defined above, with every original subdivision edge retained.

Theorem A. For every stated table and every finite host,
\fitdisplay{
EH_D^4\le e^{-c d},\qquad c=6.
}{R41.1.2}

For every independent original \(M\ge 1\) append exactly \(M\) actual \(T\) edges, \(U_D=T^M H_D\), and write \(I_q(U)=E[q^{-\frac{1}{3}}\lvert U\rvert ^{8/3}]\).

Theorem B. The SAME \(H_D\) satisfies

\fitdisplay{
I_q(U_D)^3\le
 \begin{cases}
 e^{-(2c-\frac{204}{19})d},&M=1,\\
 e^{-(2c-\frac{24427}{3420})d-5d^2},&M=2,\\
 e^{-(2c+\frac{23}{36})d},&M=3,\\
 e^{-[2c-11+\lfloor M/3\rfloor/2]d},&M\ge4.
 \end{cases}
}{R41.1.3}

In particular, every row is at most \(e^{-24d/19}\). The same
input direction applies for every independent positive \(M\), with
all equal-pair factors restored. The fully reversed statement follows
by reversing the complete configuration.

\paragraph{R41.2 Exact finite inputs, without a moment substitution}

The original C014 calculus gives

\fitdisplay{
\begin{gathered}
0\le d\le1/3,\quad d^2\le \delta_2\le\theta d,\quad\\
\delta_{i+1}\le\theta \delta_i,\quad \|B\|^2\le\theta,\quad\\
\lvert BV(x)\rvert \le\|V\|.
\end{gathered}
}{R41.2.1}

Only \(f_n,\allowbreak \allowbreak \bar f_n\) have means \(\mu_n\) in [0,1]. In particular \(E h_{2}=E W_{2}=1+d\). All four actual arm messages are nonnegative and at most two. Genuine original pairs satisfy \(f_n h_n\le 1\) and \(\bar f_n W_n\le 1\).

Let \(G_n(z)=\sum_{i=0}^{n-1} z^i,\allowbreak \allowbreak  G_0=0,\allowbreak \allowbreak  \beta_{1}=\beta_{2}=0\) and

\fitdisplay{
\beta_n=\sum_{i=1}^{\lfloor(n-1)/2\rfloor}(-1)^{i-1}\delta_i.
}{R41.L147}

The actual finite positive convolution and geometric bounds are

\fitdisplay{
\begin{gathered}
1+G_n(B)u=\bar f_n+\sum_{r=3}^n\beta_r \bar f_{n-r},\qquad\\
(1-\theta)d\le\beta_n\le d\quad(n\ge3),
\end{gathered}
}{R41.2.2}

\fitdisplay{
\begin{gathered}
\|G_n(B)\|\le1\quad(n=0,1\text{ or }n\ge3),\qquad\\
\|G_2(B)\|^2\le1+\theta=45/32.
\end{gathered}
}{R41.2.3}

The \(G_{2}\) exception is not discarded. The true operator powers satisfy

\fitdisplay{
\begin{gathered}
\|T\|\le1+32d/19,\quad \|T^2\|<\frac{7}{5},\quad\\
\|T^n\|\le(1-d/8)^{\lfloor n/3\rfloor}\quad(n\ge3).
\end{gathered}
}{R41.2.4}

The last bound uses actual blocks of lengths 3,4,5, not monotonicity of operator norms with length.

The supplied actual R2-C038/A019 terminal packets are independent of the receiving anchor. They give

\fitdisplay{
\begin{gathered}
F\le e^{R-\lambda d},\quad\\
\lambda_2=\frac{607}{160},\quad\lambda_3=\frac{767}{160},\quad\\
R=\sum_i\sum_{r=a_i+1}^{b_i}K^r\one.
\end{gathered}
}{R41.2.5}

Each mask is one finite zero-one coefficient mask. The disjoint intervals are only those within the terminal, not the union with root posts. Every extra terminal pair increases the available loss. The stated minimum losses already include \(j=2\).

For minimum 2 put \(\sigma=ER^2,\allowbreak \allowbreak  R=U(B)u,\allowbreak \allowbreak  U=\alpha+z \beta\) at \(z=i \sqrt{x}\). The actual original spectral measure \(\nu\) at \(u\) has mass \(d\), support \([0,\allowbreak \allowbreak \theta]\), and \(\int  x^r=\delta_{r+1}\). The finite packet is

\fitdisplay{
\sigma<d/3,\quad \lvert R\rvert <\frac{3}{5},\quad e^R\le1+R+5R^2/8,
}{R41.L186}

\fitdisplay{
x/2\le\alpha\le6x/5,\qquad -6x/5\le\beta\le x/2.
}{R41.2.6}

For \(D_i=E B^iR,\allowbreak \allowbreak  a_0=-D_1\) and \(D=D_0+5\sigma/4\), put \(\epsilon=D+D_1\). Then

\fitdisplay{
0\le a_0\le d/2,\quad -d/5\le D\le11d/12,\quad
 \lvert \epsilon\rvert \le7d/10,
}{R41.L195}

\fitdisplay{
\begin{gathered}
-d/5\le D_2\le d/12,\quad 0\le D_3\le d/5,\quad\\
\lvert D_i\rvert \le d/12\ (i\ge4),\quad D_0\le d/2.
\end{gathered}
}{R41.2.7}

These are actual joint quantities, not independently chosen parameters. The positive upper and complete finite losses are

\fitdisplay{
\begin{gathered}
e^{\lambda_2d}T^LF\le Y_L\\
=\bar f_L+B^LR+D \bar f_{L-1}+\sum_{i=1}^{L-1}D_i \bar f_{L-1-i},
\end{gathered}
}{R41.2.8}

\fitdisplay{
\begin{gathered}
Y_1=q+D+BR,\quad\\
Y_L=1+\epsilon+Q_L(B)u-\sum_{r=3}^L\gamma_r \bar f_{L-r}\quad(L\ge2),
\end{gathered}
}{R41.L211}

\fitdisplay{
\begin{gathered}
Q_n=G_n+z^nU+\epsilon G_{n-2}+Dz^{n-2}\quad(n\ge2),\quad\\
Q_{n+1}=1+\epsilon+zQ_n,\quad E Q_n(B)u=-\gamma_{n+1},
\end{gathered}
}{R41.2.9}

\fitdisplay{
\gamma_n=\beta_n+D\beta_{n-1}-a_0\beta_{n-2}-D_{n-1}\quad(n\ge3).
}{R41.2.10}

Thus \(\gamma_{3}=d-D_2,\allowbreak \allowbreak  \gamma_{4}=d+Dd-D_3\). The actual uniform bounds are

\fitdisplay{
\begin{gathered}
11d/12\le\gamma_3\le6d/5,\quad\\
\gamma_4\ge11d/15,\quad\\
d/4\le\gamma_n\le25d/18\quad(n\ge4).
\end{gathered}
}{R41.2.11}

The last upper includes \(\gamma_{4}\), since \(\gamma_{4}\le d+11d^2/12\). For \(n\ge 5\) it follows from \(\beta\le d,\allowbreak \allowbreak  D\le 11d/12\) and \(\lvert D_{n-1}\rvert \le d/12\). Every lower bound retains \(\beta_{1}=\beta_{2}=0\) at its actual short endpoint.

For minimum 3 the supplied packet gives

\fitdisplay{
\sigma<9d/64,\ \lvert R\rvert <\frac{1}{3},\quad e^R\le1+R+3R^2/5,
}{R41.L236}

\fitdisplay{
D_0\in [-d/5,d/2],\quad D_1\in [-d/12,d/5],\quad
 -D_2\in [-d/12,d/5].
}{R41.2.12}

For \(S_L=T^L\) or \(A T^{L-1}\), the actual positive upper is

\fitdisplay{
\begin{gathered}
0\le S_LF\le e^{-\lambda_3d}\\
(S_L1+S_LR+\omega_L),\quad\\
\omega_1<27d/160,\quad\omega_{L\ge2}<27d/80.
\end{gathered}
}{R41.2.13}

Here \(E AR\le 7d/12,\allowbreak \allowbreak  E ATR\le 9d/10\), and the TRUE propagated square gives

\fitdisplay{
\begin{gathered}
TR=D_0\,1+BR,\quad\\
\|TR\|^2=D_0^2+2D_0D_1+\|BR\|^2<(\frac{11}{24})^2d,\\
\quad \|T^{L-1}R\|<(\frac{21}{40})\sqrt d\quad(L\ge3).
\end{gathered}
}{R41.2.14}

The genuine constant cross in \(TR\) is part of this bound.

All negative root posts have one actual finite envelope

\fitdisplay{
\begin{gathered}
P\le e^{S-\Gamma_r},\quad\\
\Gamma_r\ge16kd/5,\quad S=BV_p,\quad\\
ES<d/12,\quad \|V_p\|<3\sqrt d/8\quad(k\ge1).
\end{gathered}
}{R41.2.15}

It follows from the actual finite negative pair mask \(B^4 Q_p(B)u\). The squared norm of \(V_p\) is at most \((\frac{2197}{16245})d\), and its mean is bounded by \((\frac{1024}{855})\delta_4<d/12\). Overlap with the terminal mask is allowed and is paid below by adding the vectors.

\paragraph{R41.3 A separate fourth-exponent payment}

The supplied all-real signed lemma gives, for EVERY original real \(V\),

\fitdisplay{
E e^{4BV}\le\exp\{4EBV+(\frac{17}{4})\|V\|^2\}.
}{R41.3.1}

It includes actual constant coordinates and either sign of \(EBV\). Suppose \(X\) is the actual nonnegative anchor and \(S_LF\le e^{-\lambda d}Y\) with \(Y\ge 0\), and the actual sum has upper

\fitdisplay{
X+Y\le2+c_0+BV.
}{R41.L282}

AM--GM and \(1+t\le e^t\) applied to the nonnegative average imply

\fitdisplay{
\log E(X S_LF)^4
 \le-4\lambda d+4(c_0+EBV)+(\frac{17}{4})\|V\|^2.
}{R41.3.2}

For posts, the SAME calculation gives

\fitdisplay{
\begin{gathered}
\log E(PX S_LF)^4\le\\
-4\lambda d-4\Gamma_r+4(c_0+EBV)+4ES\\
+(\frac{17}{4})\|V+V_p\|^2.
\end{gathered}
}{R41.3.3}

 Genuine equal pairs can then be restored by positivity. These formulas, not an eighth-to-fourth upgrade, prove all new fourth bounds.

\paragraph{R41.4 All \(m\ge 3\): complete fourth payment for the R2-C038 directions}

The positive sums in R2-C038/A019 give the following fourth-moment estimates. The coefficients track the change of exponent. All entries are in units of \(d\). Minimum 2 starts with \(4\lambda_2=\frac{607}{40}\).


\begingroup\small
\setlength{\tabcolsep}{3pt}
\setlength{\LTleft}{0pt}\setlength{\LTright}{\fill}
\begin{longtable}{@{}p{0.2200\linewidth}p{0.1750\linewidth}p{0.1750\linewidth}p{0.1750\linewidth}p{0.1750\linewidth}@{}}
\caption{Short even lists: the fourth moment and all outer lengths: parameter values.}\label{tab:parameters:r2-c041:2}\\
\toprule
Entry & Bare fourth loss & Other scalar fourth fee & Full squared-norm cost & Resulting loss \\
\midrule\endfirsthead
\toprule
Entry & Bare fourth loss & Other scalar fourth fee & Full squared-norm cost & Resulting loss \\
\midrule\endhead
\bottomrule\endfoot
\(L=1,\allowbreak \allowbreak m=3\) & 8; & \(\frac{17}{3}\); & \(<\frac{15}{2}\); & \(>\frac{1201}{120}\). \\
\addlinespace[2pt]
\(L=1,\allowbreak \allowbreak m\ge 4\) & \(\frac{19}{4}\); & \(\frac{17}{3}\); & \(\frac{51}{8}\); & \(\frac{473}{60}\). \\
\addlinespace[2pt]
\(L=2\) & \(\frac{1}{5}\); & 0; & \(<\frac{53}{8}\); & \(>\frac{35}{4}\). \\
\addlinespace[2pt]
\(L=3\) & \(\frac{47}{15}\); & 0; & \(<\frac{29}{4}\); & \(>\frac{1327}{120}\). \\
\end{longtable}
\endgroup


\(At L=1\) the actual vectors are \(G_{m-1}u-R\). Their squared norms are \(\le 167d/96\) for \(m=3\) and \(\le 3d/2\) for \(m\ge 4\); the full signed scalar fee is \(4(\omega_1+EAR)\le 17d/3\). \(At L=2\) the true vector is \(D\one-B(G_{m-2}u-R)\), with squared norm \(\le 3329d/2160\) and total signed mean \(\le -d/20\). \(At L=3\) the true vector is \((G_{2}-G_{m-1})u+D u+B^2R+\epsilon\one\), with squared norm \(<169d/100\) and total signed mean \(\le -47d/60\). Thus every table entry follows directly from (R41.3.2) with its actual constant cross.

For every \(L\ge 4\) put \(t=\epsilon/d-\frac{1}{4}\) in \([-\frac{19}{20},\allowbreak \allowbreak \frac{9}{20}]\) and \(C=\frac{113}{100}\). The full actual mean and norm satisfy

\fitdisplay{
(c_0+EBV)/d\le t-\frac{1}{3},\quad
 \|V\|^2/d\le N(t)=
 \begin{cases}
 C^2+(t^2-2t)/3,&t\le0,\\
 C^2+(t^2+3t/2)/3,&t\ge0.
 \end{cases}
}{R41.4.1}

These follow from the exact \(Q_{L-1}-G_{m-1}\) vector, whose norm is \(<(\frac{1}{2}+\frac{3}{20}+\frac{7}{25}+\frac{1}{5})\sqrt{d}\), and whose mean lies \([-d,\allowbreak \allowbreak 3d/4]\). The fourth fee \(4(t-\frac{1}{3})+(\frac{17}{4})N(t)\) is convex on each whole half. \(At\) the negative endpoint, zero and positive endpoint respectively,

\fitdisplay{
\begin{gathered}
(\frac{17}{4})C^2<\frac{87}{16},\quad\\
(\frac{17}{4})(\frac{1121}{1200})<4,\quad\\
(\frac{17}{4})(\frac{117}{400})<5/4.
\end{gathered}
}{R41.L340}

The fourth mean fees at those endpoints are \(-\frac{77}{15},\allowbreak \allowbreak -\frac{4}{3},\allowbreak \allowbreak \frac{7}{15}\). The resulting fees are respectively \(<\frac{71}{16},\allowbreak \allowbreak <5,\allowbreak \allowbreak <\frac{115}{16}\). Subtracting from \(\frac{607}{40}\) leaves more than 6 in every case. Convexity proves every actual \(\epsilon\).

For minimum 3 the actual R2-C038 bare sums and propagated norms give:


\begingroup\small
\setlength{\tabcolsep}{3pt}
\setlength{\LTleft}{0pt}\setlength{\LTright}{\fill}
\begin{longtable}{@{}p{0.2200\linewidth}p{0.1750\linewidth}p{0.1750\linewidth}p{0.1750\linewidth}p{0.1750\linewidth}@{}}
\caption{Short even lists: the fourth moment and all outer lengths: parameter values.}\label{tab:parameters:r2-c041:3}\\
\toprule
Entry & Start including bare loss & Norm cost & Other scalar fee & Net loss \\
\midrule\endfirsthead
\toprule
Entry & Start including bare loss & Norm cost & Other scalar fee & Net loss \\
\midrule\endhead
\bottomrule\endfoot
\(L=1,\allowbreak \allowbreak m=3\) & \(\frac{767}{40}+8\); & \(<\frac{43}{4}\); & \(<\frac{7}{2}\); & \(>\frac{517}{40}\). \\
\addlinespace[2pt]
\(L=1,\allowbreak \allowbreak m\ge 4\) & \(\frac{767}{40}+\frac{19}{4}\); & \(<\frac{33}{4}\); & \(<\frac{7}{2}\); & \(>\frac{487}{40}\). \\
\addlinespace[2pt]
\(L=2\) & \(\frac{767}{40}+\frac{3}{4}\); & \(<\frac{27}{4}\); & \(<5\); & \(>\frac{327}{40}\). \\
\addlinespace[2pt]
\(L\ge 3\) & \(\frac{767}{40}+\frac{19}{4}\); & \(<\frac{9}{2}\); & \(<10\); & \(>\frac{377}{40}\). \\
\end{longtable}
\endgroup


The respective vector norms are \(\frac{63}{40},\allowbreak \allowbreak \frac{11}{8},\allowbreak \allowbreak \frac{151}{120},\allowbreak \allowbreak \frac{41}{40}\) times \(\sqrt{d}\). The last row uses the actual \(W_{\min(m,\allowbreak \allowbreak L)}+\bar f_{\max(m,\allowbreak \allowbreak L)}\) sum. These rows include equality \(m=L\) and the \(G_{2}\) and \(T^2\) exceptions.

For \(k\ge 1\) posts add \(\frac{64}{5}\) to the starting fourth loss and at most \(\frac{1}{3}\) to the mean fee. In minimum 2's \(L=1,\allowbreak \allowbreak L=2,\allowbreak \allowbreak L=3\) rows the full vector norms are respectively \(<\frac{41}{24},\allowbreak \allowbreak \frac{13}{8},\allowbreak \allowbreak \frac{67}{40}\) times \(\sqrt{d}\); their costs are \(<\frac{25}{2},\allowbreak \allowbreak <\frac{45}{4},\allowbreak \allowbreak <12\). The resulting losses exceed \(\frac{569}{40},\allowbreak \allowbreak 6,\allowbreak \allowbreak \frac{751}{40}\). The \(L=2\) explicit lower bound is \(\frac{1119}{40}+\frac{1}{5}-\frac{45}{4}-\frac{1}{3}>6\).

For minimum \(2 L\ge 4\), complete the squares in (R41.4.1):

\fitdisplay{
\begin{gathered}
N=C^2-\frac{1}{3}+(t-1)^2/3\ (t\le0),\quad\\
N=C^2-\frac{3}{16}+(t+\frac{3}{4})^2/3\ (t\ge0).
\end{gathered}
}{R41.L369}

Both residual constants are positive. Hence \(\sqrt{N}\), and then \((\sqrt{N}+\frac{3}{8})^2\), are convex on each half. \(At\) the negative, zero, positive endpoints the fourth norm costs are \(<15,\allowbreak \allowbreak <\frac{39}{4},\allowbreak \allowbreak <\frac{23}{2}\). Including means and \(ES\), the full fees are \(<\frac{51}{5},\allowbreak \allowbreak <\frac{35}{4},\allowbreak \allowbreak <\frac{123}{10}\). The starting loss is \(\frac{1119}{40}\), so all leave more than 6.

For minimum 3 posts the full vector norms are \(\frac{39}{20},\allowbreak \allowbreak \frac{7}{4},\allowbreak \allowbreak \frac{49}{30},\allowbreak \allowbreak \frac{7}{5}\) in the same four rows. The starting losses including bare terms are \(\frac{1599}{40},\allowbreak \allowbreak \frac{1469}{40},\allowbreak \allowbreak \frac{1309}{40},\allowbreak \allowbreak \frac{1469}{40}\); norm costs are respectively \(<\frac{65}{4},\allowbreak \allowbreak <\frac{53}{4},\allowbreak \allowbreak <\frac{23}{2},\allowbreak \allowbreak <\frac{17}{2}\), and other fees \(<\frac{7}{2},\allowbreak \allowbreak <\frac{7}{2},\allowbreak \allowbreak <\frac{11}{2},\allowbreak \allowbreak <\frac{21}{2}\). Their differences all exceed 6. Extra pairs only improve \(\Gamma\). This proves (R41.1.2) for every \(m\ge3\).

\paragraph{R41.8 Every independent outer length from the same fourth input}

For any original positive kernel operator \(P\), put \(h=P_{1}\). Original row Cauchy and \(H\)ölder with exponents 3 and \(\frac{3}{2}\) give, for every real \(Y\),

\fitdisplay{
\begin{gathered}
\lvert PY\rvert ^2\le hP(\lvert Y\rvert ^2),\quad\\
I_q(PY)^3\le E(h^4/q)\|P\|^4(E\lvert Y\rvert ^4)^2.
\end{gathered}
}{R41.8.1}

All quantities remain on \(L=2(\mu)\), with actual adjoint \(T^*=A\).

\(At M=1,\allowbreak \allowbreak  h=q\). Since \(Eu=0,\allowbreak \allowbreak  Eu^2=d\) and \(\lvert u\rvert \le 1\),

\fitdisplay{
E q^3=1+3d+Eu^3\le1+4d.
}{R41.L671}

Together with (R41.2.4), this pays the complete cost \(4+\frac{128}{19}=\frac{204}{19}\) in (R41.8.1). The ordinary-fourth bridge combines the displayed row and norm costs
with the unweighted fourth moment of the input. Its input bound is
distinct from the \(p-\) or \(q-\) weighted moments.

For \(M=2\), the exact C094 whole-host calculation gives

\fitdisplay{
I_q(T^2Y)^3\le
 e^{(\frac{24427}{3420})d-5d^2}(E\lvert Y\rvert ^4)^2.
}{R41.8.2}

Its negative curvature combines the \(\bar f_{2}\) inverse-fourth row with the same \(\delta_2\).

For \(M=3\), C094, Section 2.3 gives the actual row \(E(\bar f_{3}^4/q)\le 1-23d/36\) and \(\lVert T^3\rVert \le 1\). Hence

\fitdisplay{
I_q(T^3Y)^3\le(1-23d/36)(E\lvert Y\rvert ^4)^2.
}{R41.8.3}

For every \(M\ge 4\), the exact C084 row used in the same C094 section is \(E(\bar f_M^4/q)\le e^{11d}\). The fourth power of the true operator bound retains the floor:

\fitdisplay{
I_q(T^M Y)^3\le
 e^{[11-\lfloor M/3\rfloor/2]d}(E\lvert Y\rvert ^4)^2.
}{R41.8.4}

The \(M=3\) and \(M\ge 4\) rows are distinct.

Substitute the actual \(H_D\) of Theorem A into
(R41.8.1)--(R41.8.4). For \(c=6\), the losses are
\(24/19\) at \(M=1\), \(16613/3420\) with the additional term
\(5d^2\) at \(M=2\), \(455/36\) at \(M=3\), and
\(1+\lfloor M/3\rfloor/2\ge3/2\) at \(M\ge4\).
These are the bounds in (R41.1.3), each at least \(24/19\).

